\section{模型假设}

\begin{enumerate}
    \item 两路传感器观测同一真实轨迹 $\boldsymbol{p}(t)$，坐标系之间无旋转分量，仅含平移偏差。

    \textbf{依据}：题目明确两路数据针对同一轨迹，实测数据经平移校正后轨迹形态相关系数 $> 0.999$，引入旋转参数对残差无显著改善。

    \item 以方式1为参考系（偏差零点），方式2存在时间偏差 $\Delta\tau$ 和可能的空间偏差 $(\Delta x, \Delta y)$。

    \textbf{依据}：需固定一路避免系统欠定。方式1虽采样率低但噪声水平较稳定，选为基准不失一般性。

    \item 附件2/3中观测噪声服从零均值高斯分布，各时刻独立，协方差为 $\sigma_i^2 \boldsymbol{I}$。

    \textbf{依据}：二阶差分残差直方图近似正态，时序自相关系数 $< 0.1$，XY 分量方差比接近 1:1。

    \item 机器人短时运动满足常加速度模型（0.1\,s 内加速度视为恒定）。

    \textbf{依据}：采样间隔内加速度变化量远小于位置误差量级，对状态预测的近似误差可忽略。

    \item 任务执行为瞬时动作，准备阶段（射击 1.5\,s / 拍照 0.5\,s）内约束需持续满足。

    \textbf{依据}：题目原文规定准备时长，物理上瞄准/对焦过程需全程稳定。

    \item 单机器人不可同时执行两个任务，任务时间窗不得重叠。

    \textbf{依据}：单执行器物理约束，准备阶段占用唯一控制通道。

    \item 附件数据采样时钟稳定，等间隔假设成立。

    \textbf{依据}：实测相邻时间间隔标准差 $< 10^{-13}$\,s，达浮点精度极限。
\end{enumerate}
